\section{Descripción del Servicio y Organización}
\label{sec:descripcion}

\subsection{Servicios de Cafetería}
Se establecen dos modalidades de servicio de cafetería diferenciadas para atender de manera óptima las necesidades de cada colectivo:

\subsubsection{Cafetería de Público}
Destinada a visitantes, familiares y pacientes ambulatorios. Ofrecerá una amplia gama de productos saludables, menús diarios y opciones especiales (sin gluten, vegetarianas, etc.).

\subsubsection{Cafetería de Profesionales}
Destinada al personal del centro hospitalario. Contará con tarifas bonificadas según pliegos y un sistema de servicio ágil para adaptarse a los tiempos de descanso del personal sanitario.

\subsection{Servicio de Manutención para Personal Autorizado}
El servicio de manutención está especialmente diseñado para el personal de guardia médica u otros profesionales autorizados del \textbf{Hospital de Rehabilitación y Traumatología}. 
\begin{itemize}
    \item \textbf{Desayuno y Cafés:} Dispensación en horarios establecidos directamente en las áreas habilitadas.
    \item \textbf{Almuerzo y Cena:} Menús equilibrados nutricionalmente, supervisados por el servicio de dietética del hospital.
    \item \textbf{Distribución en Unidades:} Protocolo de entrega de comidas en bandejas térmicas para el personal que no pueda abandonar su puesto de guardia.
\end{itemize}

\subsection{Plan de Control de Higiene y Alérgenos}
Toda la manipulación de alimentos se regirá bajo un estricto sistema de \textbf{Análisis de Peligros y Puntos de Control Crítico (APPCC)}, garantizando la trazabilidad y seguridad alimentaria en todo momento. Se prestará especial atención a la identificación de alérgenos según el Reglamento (UE) N.º 1169/2011 \cite{reglamento_alergenos}.
